\begin{textsample}{POS Dim 3 – al – Score 25.00 – al/al\_inf\_15.txt}  \label{ex:f3_pos_212}
Leitura e Escrita na Educação \textbf{Infantil} Segundo a BNCC, as Interações e a \textbf{Brincadeira} são eixos estruturantes que promovem aprendizagens essenciais, compreendendo comportamentos, habilidades, conhecimentos e vivências que se estabelecem enquanto aprendizagem e desenvolvimento nos diversos campos de experiências a partir do pleno exercício dos direitos de aprendizagens presentes no cotidiano das práticas pedagógicas, incluindo nessas práticas, as vivências com a leitura e a escrita de forma singular ao desenvolvimento \textbf{infantil} presente nessa Etapa da Educação Básica.
Novas expectativas nascem em relação ao trabalho com a leitura e a escrita frente às iniciativas educativas normatizadas pela BNCC para atender as demandas da sociedade atual.
Seja através de imagens, \textbf{sons}, falas e escritas, explorando distintas linguagens como forma de expressão da \textbf{criança}, essa etapa da Educação Básica move -se no desenvolvimento \textbf{infantil} provocando \textbf{curiosidades} e instigando interesses.
Dessa forma, propomos nesse referencial possibilidades de experiências com a linguagem materna – através de vivencias sociais de leitura, escrita e oralidade – de modo que habilidades, noções, valores e \textbf{afetos} sejam trazidos para os espaços e ambientes de convivências sem queimar etapas, respeitando as múltiplas \textbf{infâncias} que singularizam e definem as \textbf{crianças} de nosso Estado.
É nesse movimento, que os arranjos curriculares presentes na Base promovem um nortear de descobertas sobre a língua materna e sobre os gêneros literários, textos orais e escritos existentes em nossa cultura alagoana, possibilitando para nossas \textbf{crianças} construir sua identidade num processo de cidadania e valorização da cultura local.
Compreendendo essa \textbf{criança} como sujeito histórico e de direitos, identificamos nos acervos literários e culturais de nosso Estado, e também de nosso Brasil, como o lugar de \textbf{imersão} cultural que compreende a \textbf{criança} de zero a cinco anos de idade enquanto sujeito de linguagem.
Guinchar a \textbf{criança} para esse lugar de falante, escritor ( escrita espontânea ) e leitor, é considerá -la enquanto pessoa que se comunica, se expressa, vive, convive, (re)cria relações, (re)significando a formação do “ Eu ”, do seu “ EU ” autônomo que se (re)conhece num falar próprio, que não foge do padrão, e sim o enriquece com a melodia do falar alagoano, com a riqueza das expressões nordestinas, onde a linguagem e os acervos e gêneros literários só favorecem uma relação com a leitura e escrita repleta de descobertas e \textbf{brincadeiras}.
A relação de nossas \textbf{crianças} a partir de \textbf{bebês} com a língua se dá \textbf{cheia} de tradições e expressões próprias do nosso território, onde o diminutivo é expressão de carinho, como por exemplo : mainha e painho.
E nesse aconchego a linguagem trás para os espaços e ambientes de vivências experiências com palavras \textbf{cheias} de significados, como identificar -se com o próprio nome, como também, identificar o outro que difere do seu, formando um nós repleto de diversidade que se completa, externando a identidade de um povo, a riqueza de ser alagoano.
Desse modo, trabalhar o nome das \textbf{crianças} em Alagoas através de fichas, jogos, músicas, danças e \textbf{brincadeiras}, é trazer a identidade da família inteira – de geração em geração, é evidenciar as raízes quilombolas e/ou indígenas, que mesmo que não se veja na cor da pele, no liso ou crespo do cabelo, essa identidade está no sangue quando, por exemplo, colocamos as \textbf{crianças} diante de uma \textbf{roda} de Toré e elas dançam com todo entusiasmo, coerência e movimentos precisos.
Todo trabalho realizado nos centros de Educação \textbf{Infantil} e Escolas de nosso Estado, precisam se apropriar desse trabalho com a linguagem oral e escrita que faça sentido para \textbf{criança} enquanto ser cultural e produtora de cultura, sendo para elas experiências significativas.
De maneira evolutiva, em suas tentativas, as \textbf{crianças} precisam passar a fazer um uso mais complexo da linguagem frente a cada nova fase de desenvolvimento \textbf{infantil}, assim : [... ] da utilização de poucas palavras para frases, de assuntos concretos para outros mais abstratos, de situações contextualizadas no presente para situações do passado e do futuro.
Essa progressão se dá a partir das interações comunicativas de qualidade e positivas que as \textbf{crianças} têm a oportunidade de vivenciar em seu cotidiano.
Nesse contexto, é muito importante que as \textbf{crianças} bem \textbf{pequenas} tenham diferentes oportunidade de interagir com outras \textbf{crianças} e demais pessoas, falando sobre suas experiências pessoais, relatando fatos significativos, sendo \textbf{escutadas} e \textbf{acolhidas} naquilo que comunicam, expressando -se e comunicando -se por meio do corpo, do movimento, da dança, da \textbf{mímica}, do \textbf{som}, da música, de suas esculturas, desenhos ou do teatro. ( BNCC comentada, campo de experiência \textbf{Escuta}, fala pensamento e imaginação ) Por isso, é tão importante que seja criado, nos ambientes educativos, um clima seguro e engajador para que se expressem livremente, e que os \textbf{adultos} que compõem esse espaço estejam disponíveis para \textbf{conversar} e interagir com elas, sendo responsivos às suas colocações e criando um efetivo diálogo, fazendo circular a linguagem de forma \textbf{interessante}, prazerosa e instigante em todas as suas expressões e manifestações.
Para engajar essa \textbf{criança} da primeira \textbf{infância} em (re)conhecer o mundo de fantasia da Literatura e de acervos culturais de sua história não é necessário que ela saiba ler e escrever convencionalmente.
Pois para encantar uma \textbf{criança} com \textbf{livros}, textos de diversos gêneros literários e historias \textbf{contadas} de geração em geração, basta possibilitar um ambiente de experimentação.
Uma \textbf{criança} é capaz de folhear um \textbf{livro}, ver as figuras e imaginar uma história, existindo quem a \textbf{escute}, ela pode até \textbf{contar} essa história.
Talvez, sua narrativa seja diferente da que está escrita, ou talvez seja praticamente igual, caso tenha vivenciado momentos de \textbf{escuta}, ao ponto de memorizar fragmentos das histórias \textbf{ouvidas}.
Ela pode até não está lendo dentro dos padrões do sistema notacional, mas está se envolvendo e construindo sentido.
Dessa forma, evidenciamos o quão importante é na Educação \textbf{Infantil} a leitura e escrita dentro de um trabalho com literatura \textbf{infantil} e acervos relacionados à cultura local numa perspectiva de \textbf{brincadeira} e descobertas, em que as \textbf{crianças} experimentam e vivenciam o ler e escrever além das letras, dos signos e formas, (re)construindo hipóteses da funcionalidade do que está escrito e para que e/ou quem escrever.
Nesse contexto, disponibilizar \textbf{livros} para as \textbf{crianças} \textbf{manusearem}, expor as \textbf{crianças} a momentos de contação de histórias clássicas e do contexto cultural, possibilitar interações e \textbf{brincadeiras} com repentes, músicas e danças da sua localidade e de outras regiões, poemas, parlendas, travas-línguas, fazer visitas a \textbf{bibliotecas} e salas de leitura para que as \textbf{crianças} escolham os portadores de texto que quiserem, folheando, mostram para os \textbf{colegas} ( momento de livre exploração ), solicitando ao adulto-leitor para que leia, envolva as \textbf{crianças} e favoreça um pensar sobre a leitura e a escrita, formulando hipóteses e construindo noções, valores e habilidades.

% matched lemmas: acolher, adulto, afeto, bebê, biblioteca, brincadeira, cheio, colega, contar, conversar, criança, curiosidade, escuta, escutar, imersão, infantil, infância, interessante, livro, manusear, mímico, ouvir, pequeno, roda, som
\end{textsample}
